% GENERATED FILE. Edit canonical lessons, then run npm run generate:books.
% canonical-source-hash: fnv1a64:e88c5388513a9619
% canonical-lessons: GU-C48-knee, GU-C48-heart, GU-C48-blood, GU-C48-fever, GU-C48-cough

\chapter{Knee, Heart and a Fever}
\label{ch:gu-knee-heart-and-a-fever}
\hlchaptermodality{48}

\begin{chapteropening}
\textbf{By the end of this chapter:} \emph{I can name the knee, the heart and blood, and say I have a fever or a cough.}

Complete the last lesson of chapter 48: I can name the knee, the heart and blood, and say I have a fever or a cough.
\end{chapteropening}

\section[\texorpdfstring{\gu{ઘૂંટણ} (ghū̃ṭaṇ)}{ઘૂંટણ (ghū̃ṭaṇ)}]{\gu{ઘૂંટણ} (ghū̃ṭaṇ) — a knee}

\label{lesson:GU-C48-knee}



\begin{quote}
Before the new one: say the Gujarati for a lip, then the Gujarati for the tongue.
\end{quote}

\subsection*{You'll want to know: \gu{ઘૂંટણ}}
\textbf{\gu{ઘૂંટણ}} (\emph{ghū̃ṭaṇ}) — \textquotedblleft{}a knee\textquotedblright{}. Halfway down the \textbf{\gu{પગ}}.

\begin{cousinweb}
\textbf{\gu{ઘૂંટણ}} (\emph{ghūṁṭaṇ}) is a tadbhava. CDIAL lists it under \emph{ghuṇṭa}, a variant of the lexical Sanskrit \emph{ghuṭa}, \textquotedblleft{}ankle\textquotedblright{} (CDIAL 4479). It takes the word from Dravidian, after T. Burrow.
\end{cousinweb}

\begin{grammarlens}[title={What you've built}]
One more word for knee, heart and a fever.
\end{grammarlens}

\subsection*{Your turn}
\begin{itemize}
\raggedright
  \item \emph{Say it:} \emph{ghū̃ṭaṇ}
  \item \emph{Say it:} \emph{ghū̃ṭaṇ}, once more
  \item \emph{Say it:} say \emph{ghū̃ṭaṇ}
  \item \emph{Recall:} say the Gujarati for a lip, then the Gujarati for the tongue, then say \emph{ghū̃ṭaṇ} again
\end{itemize}

\begin{tcolorbox}[breakable,colback=teal!4,colframe=teal!35!black,title={Before you move on}]
What does \gu{ઘૂંટણ} mean? (\textquotedblleft{}A knee\textquotedblright{}.) Say it once more.
\end{tcolorbox}
\section[\texorpdfstring{\gu{હૃદય} (hṛday)}{હૃદય (hṛday)}]{\gu{હૃદય} (hṛday) — the heart}

\label{lesson:GU-C48-heart}



\begin{quote}
Before the new one: say the Gujarati for the tongue, then the Gujarati for a knee.
\end{quote}

\subsection*{You'll want to know: \gu{હૃદય}}
\textbf{\gu{હૃદય}} (\emph{hṛday}) — \textquotedblleft{}the heart\textquotedblright{}.

\begin{cousinweb}
\textbf{\gu{હૃદય}} (\emph{hṛday}) is Sanskrit \emph{hṛdaya}, \textquotedblleft{}heart\textquotedblright{}, taken over whole as a learned loan, a tatsama (CDIAL 14152). CDIAL lists the inherited Gujarati form as \emph{haiyũ}.
\end{cousinweb}

\begin{grammarlens}[title={What you've built}]
One more word for knee, heart and a fever.
\end{grammarlens}

\subsection*{Your turn}
\begin{itemize}
\raggedright
  \item \emph{Say it:} \emph{hṛday}
  \item \emph{Say it:} \emph{hṛday}, once more
  \item \emph{Say it:} say \emph{hṛday}
  \item \emph{Recall:} say the Gujarati for the tongue, then the Gujarati for a knee, then say \emph{hṛday} again
\end{itemize}

\begin{tcolorbox}[breakable,colback=teal!4,colframe=teal!35!black,title={Before you move on}]
What does \gu{હૃદય} mean? (\textquotedblleft{}The heart\textquotedblright{}.) Say it once more.
\end{tcolorbox}
\section[\texorpdfstring{\gu{લોહી} (lohī)}{લોહી (lohī)}]{\gu{લોહી} (lohī) — blood}

\label{lesson:GU-C48-blood}



\begin{quote}
Before the new one: say the Gujarati for a knee, then the Gujarati for the heart.
\end{quote}

\subsection*{You'll want to know: \gu{લોહી}}
\textbf{\gu{લોહી}} (\emph{lohī}) — \textquotedblleft{}blood\textquotedblright{}.

\begin{cousinweb}
\textbf{\gu{લોહી}} (\emph{lohī}) is a tadbhava, a word inherited from Sanskrit and reshaped by sound change. It comes from \emph{lohita}, \textquotedblleft{}red; blood\textquotedblright{} (CDIAL 11165). It is a form of \emph{rohita}, \textquotedblleft{}red\textquotedblright{}.
\end{cousinweb}

\begin{grammarlens}[title={What you've built}]
One more word for knee, heart and a fever.
\end{grammarlens}

\subsection*{Your turn}
\begin{itemize}
\raggedright
  \item \emph{Say it:} \emph{lohī}
  \item \emph{Say it:} \emph{lohī}, once more
  \item \emph{Say it:} say \emph{lohī}
  \item \emph{Recall:} say the Gujarati for a knee, then the Gujarati for the heart, then say \emph{lohī} again
\end{itemize}

\begin{tcolorbox}[breakable,colback=teal!4,colframe=teal!35!black,title={Before you move on}]
What does \gu{લોહી} mean? (\textquotedblleft{}Blood\textquotedblright{}.) Say it once more.
\end{tcolorbox}
\section[\texorpdfstring{\gu{તાવ} (tāv)}{તાવ (tāv)}]{\gu{તાવ} (tāv) — a fever}

\label{lesson:GU-C48-fever}



\begin{quote}
Before the new one: say the Gujarati for the heart, then the Gujarati for blood.
\end{quote}

\subsection*{You'll want to know: \gu{તાવ}}
\textbf{\gu{તાવ}} (\emph{tāv}) — \textquotedblleft{}a fever\textquotedblright{}.

\begin{cousinweb}
\textbf{\gu{તાવ}} (\emph{tāv}) is a tadbhava, a word inherited from Sanskrit and reshaped by sound change. It comes from \emph{tāpa}, \textquotedblleft{}heat, glow; pain\textquotedblright{} (CDIAL 5767).
\end{cousinweb}

\begin{grammarlens}[title={What you've built}]
One more word for knee, heart and a fever.
\end{grammarlens}

\subsection*{Your turn}
\begin{itemize}
\raggedright
  \item \emph{Say it:} \emph{tāv}
  \item \emph{Say it:} \emph{tāv}, once more
  \item \emph{Say it:} say \emph{tāv}
  \item \emph{Recall:} say the Gujarati for the heart, then the Gujarati for blood, then say \emph{tāv} again
\end{itemize}

\begin{tcolorbox}[breakable,colback=teal!4,colframe=teal!35!black,title={Before you move on}]
What does \gu{તાવ} mean? (\textquotedblleft{}A fever\textquotedblright{}.) Say it once more.
\end{tcolorbox}
\section[\texorpdfstring{\gu{ખાંસી} (khā̃sī)}{ખાંસી (khā̃sī)}]{\gu{ખાંસી} (khā̃sī) — a cough}

\label{lesson:GU-C48-cough}



\begin{quote}
Before the new one: say the Gujarati for blood, then the Gujarati for a fever.
\end{quote}

\subsection*{You'll want to know: \gu{ખાંસી}}
\textbf{\gu{ખાંસી}} (\emph{khā̃sī}) — \textquotedblleft{}a cough\textquotedblright{}. Often comes with a \textbf{\gu{તાવ}}.

\begin{cousinweb}
\textbf{\gu{ખાંસી}} (\emph{khāṁsī}) is a tadbhava of Sanskrit \emph{kāsa}, \textquotedblleft{}cough\textquotedblright{} (CDIAL 3138). CDIAL explains its \emph{kh-} as imitative, the sound of the cough itself.
\end{cousinweb}

\begin{grammarlens}[title={What you've built}]
One more word for knee, heart and a fever.
\end{grammarlens}

\subsection*{Your turn}
\begin{itemize}
\raggedright
  \item \emph{Say it:} \emph{khā̃sī}
  \item \emph{Say it:} \emph{khā̃sī}, once more
  \item \emph{Say it:} say \emph{khā̃sī}
  \item \emph{Recall:} say the Gujarati for blood, then the Gujarati for a fever, then say \emph{khā̃sī} again
\end{itemize}

\begin{tcolorbox}[breakable,colback=teal!4,colframe=teal!35!black,title={Before you move on}]
What does \gu{ખાંસી} mean? (\textquotedblleft{}A cough\textquotedblright{}.) Say it once more.
\end{tcolorbox}
